\chapter{Garis tegak lurus}\label{chap:perp}

\section{Sudut siku-siku, lancip, dan tumpul}

\begin{itemize}
\item Jika $|\measuredangle A O B|=\tfrac\pi2$, sudut $\angle A O B$ disebut \index{siku-siku!sudut}\emph{siku-siku};
%\item Jika $\measuredangle A O B\ne\pm\tfrac\pi2$, sudut $\angle A O B$ disebut \index{sudut!sudut serong}\emph{serong};
\item Jika $|\measuredangle A O B|<\tfrac\pi2$, sudut $\angle A O B$ disebut
\index{lancip}\emph{lancip};
\item Jika $|\measuredangle A O B|>\tfrac\pi2$, sudut $\angle A O B$ disebut \index{sudut!sudut siku-siku, lancip, dan tumpul}\index{tumpul}\emph{tumpul}.
\end{itemize}

\begin{wrapfigure}[2]{r}{25mm}
\vskip-25mm
\centering
\includegraphics{mppics/pic-46}
\end{wrapfigure}

Sudut siku-siku akan ditandai dengan persegi kecil seperti pada gambar.

Jika $\angle A O B$ siku-siku, kita juga dapat mengatakan bahwa $[O A)$ \index{tegak lurus}\emph{tegak lurus} terhadap $[O B)$.
Hal ini dapat ditulis sebagai \index{38@$\perp$}$[O A)\z\perp [O B)$.

Teorema~\ref{thm:straight-angle} memungkinkan kita menyebut dua garis $(O A)$ dan $(O B)$ {}\emph{tegak lurus} jika $[O A)\z\perp [O B)$; dengan demikian, kita dapat menulis $(O A)\z\perp (O B)$.

\begin{thm}{Latihan}\label{ex:acute-obtuce}
Andaikan titik $O$ terletak di antara $A$ dan $B$ dan $X\ne O$.
Tunjukkan bahwa
$\angle XOA$ lancip jika dan hanya jika
$\angle XOB$ tumpul.
\end{thm}

\section{Garis sumbu}

Andaikan $M$ adalah titik tengah ruas $[AB]$;
secara ekuivalen, $M\in(A B)$ dan $A M \z= M B$.

Garis $\ell$ yang melalui $M$ dan tegak lurus terhadap $(AB)$ disebut \index{garis sumbu}\index{tegak lurus!garis sumbu}\emph{garis sumbu} ruas~$[AB]$.

\begin{thm}[\abs]{Teorema}\label{thm:perp-bisect}
Untuk dua titik berbeda $A$ dan $B$,
semua titik yang berjarak sama dari $A$ dan $B$, dan hanya
titik-titik itu, terletak pada garis sumbu~$[A B]$.
\end{thm}

\parit{Bukti.} Misalkan $M$ adalah titik tengah~$[A B]$.
Andaikan $P A= P B$ dan $P\ne M$.
Menurut SSS (\ref{thm:SSS}),
$\triangle A M P \z\cong\triangle B M P$.
Oleh karena itu,
$\measuredangle A M P=\pm \measuredangle B M P$.
Karena $A\not=B$, tanda dalam rumus ini adalah ``$-$''.

\begin{wrapfigure}[8]{o}{34mm}
\centering
\includegraphics{mppics/pic-48}
\end{wrapfigure}

Selain itu,
\begin{align*}
\pi
&=
\measuredangle A M B
\equiv
\\
&\equiv\measuredangle A M P+\measuredangle P M B
\equiv
\\
&\equiv
2\cdot \measuredangle A M P.
\end{align*}
Dengan demikian, $\measuredangle A M P
=
\pm
\tfrac\pi2$.
Oleh karena itu, $P$ terletak pada garis sumbu.


Untuk membuktikan kebalikannya,
andaikan $P$
adalah sembarang titik pada garis sumbu $[A B]$ dan $P\z\ne M$.
Maka $\measuredangle A M P=\pm \tfrac\pi2$,
$\measuredangle B M P=\pm \tfrac\pi2$ dan
$A M\z=B M$.
Menurut SAS, $\triangle A M P\cong \triangle B M P$;
khususnya, $A P\z= B P$.
\qeds


\begin{thm}[!]{Latihan}\label{ex:pbisec-side}
Misalkan $\ell$ adalah garis sumbu $[A B]$, dan $X$ adalah sembarang titik pada bidang.
Tunjukkan bahwa
$AX<BX$ jika dan hanya jika $X$ dan $A$ terletak pada sisi yang sama dari~$\ell$.
\end{thm}

\begin{thm}[!]{Latihan}\label{ex:side-angle}
Misalkan $ABC$ adalah segitiga tak degenerat.
Tunjukkan bahwa
\[AC>BC\iff|\measuredangle ABC|>|\measuredangle CAB|.\]
\end{thm}

\begin{thm}{Latihan}\label{ex:pbisec-motion}
Suatu transformasi isometrik memiliki titik tetap $F$ dan memetakan titik $X$ ke titik lain $Y$.
Tunjukkan bahwa $F$ terletak pada garis sumbu $[XY]$.
\end{thm}




\section{Ketunggalan garis tegak lurus}

\begin{thm}[\abs]{Teorema}\label{perp:ex+un}
Terdapat tepat satu garis yang melalui suatu titik $P$ dan tegak lurus terhadap suatu garis~$\ell$.
\end{thm}

Menurut teorema di atas,
terdapat tepat satu titik $Q\in\ell$ sedemikian sehingga $(QP)\perp\ell$.
Titik $Q$ ini disebut \index{titik kaki tegak lurus}\emph{titik kaki tegak lurus} dari $P$ pada~$\ell$.

\parit{Bukti.}
Jika $P\in\ell$, keberadaan dan ketunggalan keduanya diperoleh dari Aksioma~\ref{def:birkhoff-axioms:2}.

{

\begin{wrapfigure}{r}{30mm}
\vskip-4mm
\centering
\includegraphics{mppics/pic-50}
\end{wrapfigure}

\parit{Keberadaan untuk $P\not\in\ell$.}
Misalkan $A$ dan $B$ adalah dua titik berbeda pada~$\ell$.
Pilih $P'$ sedemikian sehingga $AP'\z=AP$ dan $\measuredangle  BAP' \equiv -\measuredangle   BAP$.
Menurut Aksioma~\ref{def:birkhoff-axioms:3}, $\triangle A P' B\z\cong\triangle A P B$.
Khususnya, $A P= A P'$ dan $B P= B P'$.

Menurut Teorema~\ref{thm:perp-bisect}, $A$ dan $B$ terletak pada garis sumbu~$[P P']$.
Khususnya, $(P P')\perp (A B)=\ell$.

}


\parit{Ketunggalan untuk $P\not\in\ell$.}
Dari uraian di atas, kita dapat memilih titik $P'$ sedemikian sehingga $\ell$ merupakan garis sumbu~$[PP']$.

Andaikan $m\perp \ell$ dan $m\ni P$.
Maka $m$ merupakan garis sumbu suatu ruas $[Q Q']$ pada $\ell$;
khususnya, $P Q= P Q'$.

{

\begin{wrapfigure}{r}{37mm}
\centering
\includegraphics{mppics/pic-52}
\end{wrapfigure}

Karena $\ell$ merupakan garis sumbu $[P P']$,
kita memperoleh $PQ= P'Q$ dan $PQ' \z= P'Q'$.
Oleh karena itu,
$$P' Q=P Q=P Q'= P' Q'.$$
Menurut Teorema~\ref{thm:perp-bisect},
$P'$ terletak pada garis sumbu $[QQ']$, yaitu~$m$.
Menurut Aksioma~\ref{def:birkhoff-axioms:1}, $m=(P P')$.
\qeds

}

\begin{thm}[!]{Latihan}\label{ex:construction-perpendicular}
Uraikan konstruksi penggaris dan jangka untuk membuat garis yang melalui suatu titik dan tegak lurus terhadap suatu garis yang diberikan.
\end{thm}

\section{Pencerminan terhadap garis}

Andaikan titik $P$ dan garis $(AB)$ diberikan.
Untuk memperoleh \index{pencerminan!terhadap garis}\emph{bayangan pencerminan} $P'$ dari $P$ terhadap $(AB)$,
tarik ruas tegak lurus dari $P$ hingga mencapai $(AB)$,
lalu perpanjang ruas tersebut sejauh jarak yang sama ke sisi lainnya.

{

\begin{wrapfigure}{o}{37mm}
\vskip-2mm
\centering
\includegraphics{mppics/pic-54}
\end{wrapfigure}

Menurut Teorema~\ref{perp:ex+un}, $P'$ ditentukan secara tunggal oleh~$P$.
Perhatikan bahwa $P=P'$ jika dan hanya jika $P\in(AB)$.

\begin{thm}[\abs]{Proposisi}\label{prop:reflection}
Andaikan $P'$ adalah bayangan pencerminan $P$ terhadap $(AB)$.
Maka $AP'=AP$,
dan jika $A\ne P$,
maka
$\measuredangle BAP'\equiv -\measuredangle BAP$.
\end{thm}


\parit{Bukti.}
Jika $P\in (AB)$,
maka $P\z=P'$.
Menurut Akibat~\ref{cor:degenerate=pi}, $\measuredangle BAP\z=0$ atau~$\pi$.
Dengan demikian, pernyataan tersebut terbukti.

}

Jika $P\notin (AB)$, maka~$P'\ne P$.
Menurut konstruksi $P'$,
garis $(AB)$ merupakan garis sumbu~$[PP']$.
Oleh karena itu, menurut Teorema~\ref{thm:perp-bisect}, $AP'\z=AP$ dan $BP'\z=BP$.
Khususnya,
$\triangle ABP'\z\cong \triangle ABP$.
Dengan demikian, $\measuredangle BAP'=\pm \measuredangle BAP$.

Karena $P'\ne P$ dan $AP'=AP$,
kita memperoleh $\measuredangle BAP'\ne \measuredangle BAP$.
Jadi, satu-satunya kemungkinan yang tersisa adalah $\measuredangle BAP'\equiv-\measuredangle BAP$.
\qeds



\begin{thm}{Latihan}\label{ex:2-reflections}
Misalkan $X$ dan $Y$ masing-masing merupakan bayangan pencerminan $P$ terhadap garis $(AB)$ dan $(BC)$.
Tunjukkan bahwa $\measuredangle XBY\equiv 2\cdot \measuredangle ABC$.

\end{thm}


\begin{thm}[\abs]{Akibat}\label{cor:reflection+angle}
Pencerminan terhadap garis merupakan transformasi isometrik bidang.
Selain itu,
$\measuredangle Q'P'R'\equiv -\measuredangle QPR$
jika $\triangle P'Q'R'$ adalah bayangan pencerminan $\triangle PQR$.
\end{thm}


\parit{Bukti.}
Komposisi dua pencerminan terhadap garis yang sama
merupakan pemetaan identitas.
Khususnya, setiap pencerminan merupakan bijeksi.

Tetapkan sebuah garis $(AB)$ dan dua titik $P$ dan $Q$;
nyatakan bayangan pencerminan keduanya terhadap $(AB)$ dengan $P'$ dan $Q'$.
Mari kita tunjukkan bahwa
$$P'Q'=PQ;\eqlbl{eq:P'Q'=PQ}$$
yakni, pencerminan tersebut mempertahankan jarak.

\begin{wrapfigure}{r}{33mm}
\centering
\vskip-15mm
\includegraphics{mppics/pic-56}
\end{wrapfigure}

Tanpa mengurangi keumuman, kita dapat mengandaikan bahwa titik $P$ dan $Q$ berbeda dari $A$ dan~$B$.
Menurut Proposisi~\ref{prop:reflection}, kita memperoleh
\begin{align*}
\measuredangle BAP'&\equiv -\measuredangle BAP,
&
\measuredangle BAQ'&\equiv -\measuredangle BAQ,
\\
AP'&=AP,
&
AQ'&=AQ.
\end{align*}
Akibatnya,
\[\measuredangle P'AQ'\equiv -\measuredangle PAQ.\eqlbl{eq:P'AQ'=PAQ}\]
Menurut SAS,
$\triangle P'AQ'\cong\triangle PAQ$
dan diperoleh \ref{eq:P'Q'=PQ}.
Selain itu, kita juga memperoleh
\[\measuredangle AP'Q'\equiv\pm\measuredangle APQ.\]
Menurut \ref{eq:P'AQ'=PAQ} dan teorema tentang tanda sudut-sudut segitiga (\ref{thm:signs-of-triug}), kita memperoleh
\[\measuredangle AP'Q'\equiv-\measuredangle APQ.\eqlbl{eq:AP'Q'=APQ}\]

Dengan mengulangi argumen yang sama untuk pasangan titik $P$ dan $R$,
kita memperoleh
$$\measuredangle AP'R'\equiv-\measuredangle APR.\eqlbl{eq:AP'R'=APR}$$
Dengan mengurangkan \ref{eq:AP'R'=APR} dari \ref{eq:AP'Q'=APQ},
kita memperoleh
$\measuredangle Q'P'R'\equiv-\measuredangle QPR$.
\qedsf


\section[Isometri langsung dan tak langsung]{Transformasi isometrik langsung dan tak langsung}
\label{direct motion}

\begin{thm}[!]{Latihan kelas}\label{ex:3-reflections}
Tunjukkan bahwa setiap transformasi isometrik bidang dapat dinyatakan sebagai
komposisi paling banyak tiga pencerminan terhadap garis.
\end{thm}

Transformasi isometrik $X\mapsto X'$ disebut \index{transformasi isometrik langsung}\emph{langsung} jika
$$\measuredangle Q'P'R'= \measuredangle QPR$$
untuk setiap segitiga $PQR$;
jika sebaliknya selalu berlaku
$$\measuredangle Q'P'R'\equiv -\measuredangle QPR,$$
maka transformasi isometrik tersebut disebut \index{transformasi isometrik tak langsung}\emph{tak langsung}.

Menurut Akibat~\ref{cor:reflection+angle}, setiap pencerminan terhadap garis merupakan transformasi isometrik tak langsung.
Perhatikan bahwa komposisi dua pencerminan merupakan transformasi isometrik langsung.
Secara lebih umum, komposisi dua transformasi isometrik tak langsung bersifat langsung,
komposisi dua transformasi isometrik langsung bersifat langsung,
dan komposisi transformasi isometrik langsung dengan transformasi isometrik tak langsung bersifat tak langsung.
Khususnya, latihan di atas menyiratkan pernyataan berikut.

\begin{thm}{Proposisi}\label{prop:direct-indirect}
Setiap transformasi isometrik bidang bersifat langsung atau tak langsung.
\end{thm}


\section{Ruas tegak lurus adalah yang terpendek}
\label{sec:perp<oblique}

\begin{thm}[\abs]{Lema}\label{lem:perp<oblique}
Andaikan $Q$ adalah titik kaki tegak lurus dari $P$ pada garis~$\ell$.
Maka
$$PX>PQ$$
untuk setiap titik $X$ pada $\ell$ yang berbeda dari~$Q$.
\end{thm}

Jika $P$, $Q$, dan $\ell$ seperti di atas,
maka $PQ$ disebut \label{distance!from a point to a line}\index{jarak!dari titik ke garis}\emph{jarak dari $P$ ke~$\ell$}.

\parit{Bukti.}
Jika $P\in \ell$,
maka hasilnya langsung diperoleh karena $PQ=0$.
Jadi, kita dapat mengandaikan bahwa $P\notin \ell$.

\begin{wrapfigure}{o}{24mm}
\centering
\vskip-4mm
\includegraphics{mppics/pic-58}
\vskip0mm
\end{wrapfigure}

Misalkan $P'$ adalah bayangan pencerminan $P$ terhadap garis~$\ell$.
Perhatikan bahwa $Q$ merupakan titik tengah $[PP']$
dan $\ell$ merupakan garis sumbu $[PP']$.
Oleh karena itu,
$$PX=P'X
\quad
\text{dan}
\quad
PQ=P'Q=\tfrac12\cdot PP'.$$

Garis $\ell$ berpotongan dengan $[PP']$ hanya di titik $Q$.
Oleh karena itu, $X\notin [PP']$.
Menurut ketaksamaan segitiga dan Akibat~\ref{cor:degenerate-trig}, berlaku
$$PX+P'X>PP',$$
sehingga diperoleh hasil yang dimaksud: $PX>PQ$.
\qeds

\begin{thm}{Latihan kelas}\label{ex:right-acute}
Tunjukkan bahwa
\begin{enumerate}[(a)]
\item Jika $\triangle ABC$ memiliki sudut siku-siku di $C$, maka sudut-sudut lainnya lancip.
\item Jika $\triangle ABC$ memiliki sudut tumpul di $C$, maka sudut-sudut lainnya lancip.
\end{enumerate}
\end{thm}

Latihan ini berkaitan erat dengan Proposisi~\ref{prop:2sum}.

Jika sebuah segitiga memiliki sudut tumpul atau siku-siku, segitiga itu berturut-turut disebut \index{segitiga!segitiga siku-siku, lancip, dan tumpul}\index{tumpul}\emph{segitiga tumpul} atau \index{segitiga!segitiga siku-siku}\index{siku-siku!segitiga}\emph{segitiga siku-siku}.
Jika tidak, yaitu jika semua sudutnya lancip, segitiga itu disebut \index{lancip}\emph{segitiga lancip}.

\begin{thm}[!]{Latihan}\label{ex:obtuce}
Andaikan titik $X$ terletak di antara $V$ dan $W$.
Tunjukkan bahwa untuk setiap titik $P$, berlaku $PX < PV$ atau $PX < PW$.
\end{thm}

\begin{thm}{Latihan}\label{ex:PMQ}
Andaikan titik $P$ dan $Q$ terletak pada sisi yang sama dari garis~$\ell$.
Misalkan $M\in \ell$ adalah titik yang meminimumkan jumlah $PM+MQ$.
Tunjukkan bahwa $M\in (P'Q)$, dengan $P'$ merupakan bayangan pencerminan $P$ terhadap $\ell$.
\end{thm}



\section{Lingkaran}

Ingat bahwa lingkaran berjari-jari $r$ dan berpusat di $O$ adalah himpunan semua titik yang berjarak $r$ dari $O$.
Titik $P$ dikatakan terletak \index{di dalam!lingkaran}\emph{di dalam} lingkaran jika $OP<r$;
jika $OP>r$, titik $P$ dikatakan terletak \index{di luar lingkaran}\emph{di luar} lingkaran.
\label{def:circle}

\begin{thm}[!]{Latihan}\label{ex:inside-outside}
Misalkan $\Gamma$ adalah sebuah lingkaran dan $P\notin \Gamma$.
Andaikan sebuah garis $\ell$ melalui titik $P$
dan memotong $\Gamma$ di dua titik berbeda, $X$ dan~$Y$.
Tunjukkan bahwa $P$ berada di dalam $\Gamma$ jika dan hanya jika $P$ terletak di antara $X$ dan~$Y$.
\end{thm}

Ruas yang menghubungkan dua titik pada lingkaran disebut \index{tali busur}\emph{tali busur} lingkaran.
Tali busur yang melalui pusat lingkaran disebut \index{diameter}\emph{diameter} lingkaran.

\begin{thm}{Latihan}\label{ex:chord-perp}
Andaikan dua lingkaran berbeda $\Gamma$ dan $\Gamma'$ memiliki tali busur bersama~$[A B]$.
Tunjukkan bahwa garis yang melalui pusat $\Gamma$ dan pusat $\Gamma'$ merupakan garis sumbu~$[A B]$.
\end{thm}

\begin{thm}{Latihan}\label{ex:center}
Uraikan konstruksi penggaris dan jangka untuk memperoleh pusat
suatu lingkaran yang diberikan.
\end{thm}

{

\begin{wrapfigure}[4]{r}{54mm}
\vskip-8mm
\centering
\includegraphics{mppics/pic-60}
\end{wrapfigure}

\begin{thm}[\abs]{Lema}\label{lem:line-circle}
Sebuah garis dan sebuah lingkaran memiliki paling banyak dua titik potong.
\end{thm}


\parit{Bukti.} Andaikan $A$, $B$, dan $C$ adalah titik-titik berbeda yang terletak pada sebuah garis $\ell$ dan sebuah lingkaran $\Gamma$ berpusat di~$O$.
Maka $OA=OB=OC$; khususnya, $O$ terletak pada $m$ dan $n$,
yang masing-masing merupakan garis sumbu $[A B]$ dan $[B C]$.
Perhatikan bahwa titik tengah $[AB]$ dan titik tengah $[BC]$ berbeda.
Oleh karena itu, $m$ dan $n$ berbeda.
Hal ini bertentangan dengan ketunggalan garis tegak lurus (Teorema~\ref{perp:ex+un}).
\qeds

}

\begin{thm}[!]{Latihan}\label{ex:two-circ}
Tunjukkan bahwa dua lingkaran berbeda memiliki paling banyak dua titik potong.
\end{thm}

Sebagai akibat dari lema di atas,
sebuah garis $\ell$ dan sebuah lingkaran $\Gamma$ dapat memiliki 2, 1, atau 0 titik potong.
Dalam dua kasus pertama, garis tersebut masing-masing disebut \index{garis sekan}\emph{sekan} atau \index{garis singgung}\emph{garis singgung};
jika $P$ merupakan satu-satunya titik potong $\ell$ dan $\Gamma$,
kita mengatakan bahwa \textit{$\ell$ menyinggung $\Gamma$ di $P$}.

Serupa dengan itu, menurut Latihan~\ref{ex:two-circ},
dua lingkaran berbeda dapat memiliki 2, 1, atau 0 titik potong.
Jika $P$ merupakan satu-satunya titik potong lingkaran $\Gamma$ dan $\Gamma'$,
kita mengatakan bahwa \index{bersinggungan!lingkaran}\emph{$\Gamma$ bersinggungan dengan $\Gamma'$ di $P$}; kita juga menganggap bahwa sebuah lingkaran bersinggungan dengan dirinya sendiri di setiap titiknya.

\begin{thm}[\abs]{Lema}\label{lem:tangent}
Misalkan $\ell$ adalah sebuah garis dan $\Gamma$ adalah sebuah lingkaran berpusat di~$O$.
Andaikan $P$ merupakan titik bersama $\ell$ dan~$\Gamma$.
Maka $\ell$ menyinggung $\Gamma$ di $P$ jika dan hanya jika $(PO)\perp \ell$.
\end{thm}

\parit{Bukti.}
Misalkan $Q$ adalah titik kaki tegak lurus dari $O$ pada~$\ell$.

Andaikan~$P\ne Q$.
Misalkan $P'$ adalah bayangan pencerminan $P$ terhadap~$(OQ)$.
Perhatikan bahwa $P'\in\ell$ dan $(OQ)$ merupakan garis sumbu~$[PP']$.
Oleh karena itu, $OP=OP'$.
Dengan demikian, $P,P'\in \Gamma\cap \ell$;
artinya, $\ell$ merupakan garis sekan terhadap~$\Gamma$.

Jika $P=Q$,
maka menurut Lema~\ref{lem:perp<oblique},
$OP<OX$ untuk setiap titik $X\in \ell$ yang berbeda dari~$P$.
Dengan demikian, $P$ merupakan satu-satunya titik pada irisan $\Gamma\cap\ell$;
artinya, $\ell$ menyinggung $\Gamma$ di~$P$.
\qeds

\begin{thm}{Latihan}\label{ex:tangent-circles}
Misalkan $\Gamma$ dan $\Gamma'$ adalah dua lingkaran yang masing-masing berpusat di $O$ dan $O'$, serta berjari-jari $r$ dan~$r'$.
Tunjukkan bahwa $\Gamma$ bersinggungan dengan $\Gamma'$ jika dan hanya jika salah satu syarat berikut berlaku:
\begin{enumerate}[(a)]
\item\label{ex:tangent-circles:a} Kedua lingkaran berpotongan di sebuah titik~$P$ yang segaris dengan $O$ dan~$O'$.
\item\label{ex:tangent-circles:b} Kedua lingkaran berpotongan di sebuah titik~$P$ dan memiliki garis singgung bersama di titik ini.
\item\label{ex:tangent-circles-2} $OO' = r + r'$ atau $OO' = |r - r'|$.
\end{enumerate}

\end{thm}


{

\begin{wrapfigure}{r}{20mm}
\vskip-10mm
\centering
\includegraphics{mppics/pic-62}
\end{wrapfigure}

\begin{thm}{Latihan}\label{ex:tangent-circles-3}
Andaikan tiga lingkaran memiliki dua titik bersama.
Buktikan bahwa pusat ketiganya segaris.
\end{thm}

}

\begin{thm}{Latihan}\label{ex:tangent}
Uraikan konstruksi penggaris dan jangka untuk membuat garis-garis yang melalui suatu titik yang diberikan dan menyinggung suatu lingkaran yang diberikan.
\end{thm}

\begin{thm}{Latihan}\label{ex:tangent-circle}
Diberikan dua lingkaran $\Gamma_1$ dan $\Gamma_2$ serta sebuah ruas $[AB]$,
uraikan konstruksi penggaris dan jangka untuk membuat lingkaran berjari-jari $AB$
yang bersinggungan dengan kedua lingkaran $\Gamma_1$ dan~$\Gamma_2$.
\end{thm}
